\section{Introduction}

Fine-tuning large language models (LLMs) on cybersecurity data is a
double-edged experiment. The same instruction data that teaches a model to
explain a CVE, triage an advisory, or reason about an attack technique also
moves the model toward producing operational offensive content, up to agents
that autonomously exploit real vulnerabilities~\cite{fang2024exploit}. That
fine-tuning can \emph{degrade} a model's safety is by now well
documented~\cite{qi2023finetuning,lermen2023lora}; ElZemity et
al.~\cite{cyberllminstruct} made the cybersecurity case concrete with
CyberLLMInstruct.

That study, like most dataset construction efforts, relied on commercial,
cloud-hosted models to filter and synthesize data. This paper examines what
changes when the cloud is removed from the loop entirely. We take the
CyberLLMInstruct pipeline as our starting point and re-engineer its
LLM dependent stages to run fully on-device using Apple's MLX
framework~\cite{mlx}, then fine-tune a range of open models locally with
low-rank adaptation~\cite{lora}. Two properties of this all local setting
motivate the paper. First, no third-party service ever observes the
security data, and no provider safety filter constrains what the pipeline can
generate, removing both a confidentiality concern and a safety guardrail.
Second, the cost: the hardware is a single desktop, and the models
that make it work, including \emph{abliterated}~\cite{abliteration} variants
with their refusal behavior surgically removed, are freely downloadable - well
within reach with the majoritity of laptops from the last decade.

We deliberately push this to its uncomfortable conclusion. Alongside a 
defensive reproduction of the instruction dataset, we build an
\emph{offensive} branch (\offensivename) that strips the pipeline's
safety-alignment stage and uses an abliterated model to synthesize
exploitation write-ups over a public vulnerable code corpus, then fine-tunes
on the result. Our intent is not to ship an ``AI pentesting bot'' but to
document how cheaply and locally one can be approached, and what that implies.
This concern is not hypothetical for us: the author and other practitioners
have used these locally hosted unrestricted LLMs in confidential
red team engagements, which motivated a systematic measurement of what such
models cost to produce, in capability and in safety. We report that
measurement here; we do not report the engagements, which are confidential.
